% Set from report/draft.md section 5. The prose is the maintainer's; this file is
% its transit into LaTeX and nothing in it was rewritten.
\section{Conclusion}

\subsection{Contributions}

\textbf{A structural blindness in control layers driven by masked measurements.}
A controller whose error signal is computed as a statistic over \emph{valid}
measurements cannot see the error that would correct it, because the
measurements carrying that error are the ones validity has already discarded. In
the case measured here, a vergence controller took a robust statistic over valid
disparities in a foveal window; a target outside the search range either failed
the validity test or clamped to the window edge, so the controller had no way to
report \emph{too near} or \emph{too far}. It verged away from its target while
every measurement it could see said it was doing fine.

Nothing about this is specific to vergence, or to stereo. It arises wherever a
validity mask sits upstream of a control loop, which is a common arrangement.
The correction is not a better statistic but a channel that can report \emph{out
of range} --- something the discarded measurements are the only evidence for.

\textbf{Availability dominates selection.} Across the policy experiments, the
binary question \emph{can this be measured at all} outweighed the continuous
question \emph{how well} by 65:1 on median error. Closing the loop helped, and
its benefit was coverage over accuracy by at least 34\texttimes{} at every
occlusion level. Two objectives that never once agreed on a pixel produced
indistinguishable results. The consistent finding is that what limits an active
vision system is what it can see, not how cleverly it chooses among what it can
see.

\textbf{Geometric confinement as a confidence mechanism.} Expanding the depth
relation about the fixated disparity is exact at the fovea and degrades as the
square of the offset. That degradation term inflates the variance of any
estimate far from the fixated depth --- which is to say, precisely the
confidently-wrong matches that anti-calibrated stereo confidence produces. A
term introduced to bound approximation error turns out to function as a filter
against exactly the failure mode the confidence channel cannot report. It is a
property of the estimator, not of the sensor, and it transfers to passive stereo
unchanged.

\textbf{A synthetic stimulus is an instrument and must be characterised before
it is used to measure.} The random-dot fixture used here substituted, at every
depth discontinuity, an artefact both harder than real occlusion and able to
pass the validity test meant to catch it. Four of thirteen experiments went to
establishing this, and every result before them had to be re-read afterwards.

\subsection{What we may have measured}

The negative results are consistent and, taken at face value, they say that
uncertainty-driven active vision does not improve stereo depth estimation. We
are not confident that is the right reading, and the reason is worth stating
precisely.

\textbf{The biological inspiration transferred at the wrong level of
abstraction.} In a retina, foveation means \emph{not paying} for the periphery
--- the samples are never taken. Here it was implemented as \emph{paying, then
discounting}: every pixel matched at full resolution, then weighted down by
distance from gaze. Under that translation a foveal weight can only remove
information already bought, and the measurements show it removing. The mechanism
was not wrong. The level was.

That generalises. If sampling is free, looking is pointless. A complete
spherical capture already contains every direction; rotating the eyes about
their fixed centres returns an identical image, and there is literally nothing
to acquire. Looking matters exactly when sampling is scarce. Every experiment
reported here gave the system unlimited uniform sampling and then asked whether
choosing where to attend helped --- and under that setup the answer nearly has
to be no.

So the honest summary may not be \emph{active vision does not work}. It may be
\textbf{we posed a budget problem as an estimation problem.} The framework's
remaining untested form --- a sensor whose periphery is genuinely sampled more
coarsely than its fovea --- is not one more experiment in the series. It is the
first one that asks the question in the terms the biology poses it.

\subsection{On simulation as an instrument}

A second thread runs through this work and deserves stating on its own.

In simulation-based science the instrument is authored by the same hand as the
hypothesis. The fixture's artefact was not merely a limitation; it was
adversarial to its own detector --- self-consistent enough to pass the validity
test, wrong enough to poison the estimate, and concentrated in exactly the
region where the scientific question lived. Real data can be difficult, biased,
or insufficient. It is not written by someone who wants a particular answer.

We do not think this argues against synthetic stimuli, which offer dense ground
truth no capture provides. It argues that a synthetic stimulus deserves the same
scrutiny as an unfamiliar sensor: characterise its failure modes before trusting
a measurement made through it, and expect its worst behaviour to coincide with
the regime you care about, because that is usually the regime that is hardest to
synthesise.

\subsection{Directions}

\textbf{Variable-resolution sampling} is the priority, for the reason \S5.2
gives --- it is the only formulation in which the question is well posed. It
requires a matcher that operates on a non-uniform lattice, which is a research
problem rather than an engineering one and should be scoped as such. It also
changes the metric: the claim becomes \emph{the same posterior for less
computation}, which needs a declared computational budget that this project
never had to define.

\textbf{A controlled contrast axis at occlusion boundaries} would connect these
results to the confidence failure that motivated them, and settle whether that
failure is a property of real imagery or of high contrast --- which the two
stimulus families measured so far cannot separate.

\textbf{A moving head.} Rotations about fixed centres cannot change what a
complete sensor sees; translations can. Acquisition on an omnidirectional rig
requires the rig to move, which reintroduces a mechanical plant this project set
aside as untested rather than refuted.

\subsection{A note on what remains}

The most reusable output of this work is probably not the code, and not the
single positive result. It is a ledger of foreclosed possibilities, each with
the evidence that closed it and the cost of reopening it: the gaze objective
does not matter, verging to acquire harms and why, foveal weighting on a uniform
sensor is lossy, this stimulus cannot support occlusion claims and here is what
it substitutes instead.

Research whose deliverable is a well-ordered set of eliminated options is
awkward to write up and, we suspect, more reusable than most positive claims.
Someone building an active stereo system now knows several things not to try,
and the reasons --- which is a different and more durable kind of result than a
benchmark number.

\begin{figure}[htbp]
  \centering
  \includegraphics[width=\textwidth]{ledger.pdf}
  \caption{The thirteen foreclosures, and what later work did to each. A
  foreclosure that acquires a qualification acquires it as a new field on its
  own entry rather than as an edit to the original, so the fields are the
  timeline and this figure is built from them rather than from a hand-written
  table. Five of the thirteen were annotated later; eight have not been
  revisited. \texttt{fc-008} is the one that was \emph{qualified} --- by exp005,
  and the qualification did not survive exp007 --- and \texttt{fc-009} is the
  one whose verdict moved twice: written \textsc{additive}, corrected in place
  to \textsc{deferred} because the original overstated what had been measured,
  then answered \textsc{additive} at step 15 once a ray actually crossed the
  boundary. \emph{Derived} from the structure of \texttt{docs/state.yaml}.}
  \label{fig:ledger}
\end{figure}
